\documentclass[12pt,a4paper]{article}
\usepackage[utf8]{inputenc}
\usepackage[shortlabels]{enumitem}
\usepackage{fontawesome}
%Set page layout parameters
\usepackage[
top=4cm,
bottom=2cm,
left=2cm,
right=2cm,
headheight=2cm,
footskip=1cm,
marginparwidth=2cm]{geometry}
\usepackage{amsmath}
\usepackage{cancel}
\usepackage{hyperref}
\usepackage[dvipsnames]{xcolor}
\usepackage{multicol}
% To use color boxes
\usepackage{tcolorbox}
\usepackage{lipsum}
\tcbuselibrary{listings,theorems,skins,vignette,raster,breakable,magazine,poster,fitting,hooks}
\newcounter{ejemplo}
%Images path
\graphicspath{{images/}}
\usepackage{float} %provide us with the option H
\usepackage{fancyhdr}
\usepackage{tikz}
\usepackage{tikzpagenodes}
%to redefine figure caption 
\renewcommand{\figurename}{Figura}
%change text font
\usepackage[T1]{fontenc}
\usepackage{tgtermes} % LaTeX will use the TEX Gyre Bonum font family 


\urlstyle{same}
\usepackage{vwcol}
\definecolor{subtitle_yellow}{RGB}{252,213,129}
\definecolor{title_yellow}{RGB}{249,175,16}
\definecolor{steel_blue_title}{RGB}{93,138,168}
\definecolor{pewter_blue_subtitle}{RGB}{149,178,198}
\definecolor{ecru_circle}{RGB}{215,190,130}
\definecolor{ecru_circle1}{RGB}{68,148,173}

\newcommand{\myfirststyle}{%
	\begin{tikzpicture}[remember picture, overlay]
		
		%top border
		\fill[fill=steel_blue_title](current page.north west) rectangle ([yshift=-15mm]current page.north east);
		
		\fill[fill=pewter_blue_subtitle]([yshift=-17mm]current page.north west) rectangle ([yshift=-30mm]current page.north east);
		
		%page box
		%\draw[lightgray, fill=Rhodamine, fill opacity=0.05, line width=1mm, rounded corners=3mm] ([xshift=-5mm,yshift=5mm]current page text area.north west) rectangle ([xshift=5mm,yshift=-5mm]current page text area.south east);
		
		%bottom box
		\fill[fill=pewter_blue_subtitle](current page.south west) rectangle ([yshift=10mm]current page.south east);
		
		%page circle
		\draw[white,fill=ecru_circle1,line width=2mm] ([xshift=-5mm,yshift=25mm]current page text area.north east) circle (1); 
		
		
		%page number text        
		%\node[anchor=center] at ([xshift=-5mm,yshift=25mm]current page text area.north east) {\Huge\bfseries\sffamily\color{white}{\thepage}};
		
		%title text
		\node[anchor=west] at ([yshift=32mm]current page text area.north west) {\Huge\bfseries\sffamily\color{white}{Ley de Ohm en circuitos resistivos}};
		
		%subtitle text
		\node[anchor=west] at ([xshift=10mm,yshift=16mm]current page text area.north west) {\Large\bfseries\sffamily\color{white}{}};
		
		%footer text
		\node[anchor=center] at ([yshift=-15mm]current page text area.south) {\large\bfseries\sffamily\color{white}{Created by: Belmont Dubois}};
		
	\end{tikzpicture}
}


%create page style
\fancypagestyle{MyFirstStyle}{%
	\fancyhf{}
	\fancyhead[C]{\myfirststyle}
	\renewcommand{\headrulewidth}{0pt}
}%

\pagestyle{MyFirstStyle}

\colorlet{colexam}{red!75!black}
\newtcolorbox[use counter=ejemplo]{myexample}{%
	empty,title={Ejemplo \thetcbcounter},attach boxed title to top left,
	boxed title style={empty,size=minimal,toprule=2pt,top=4pt,
		overlay={\draw[colexam,line width=2pt]
			([yshift=-1pt]frame.north west)--([yshift=-1pt]frame.north east);}},
	coltitle=colexam,fonttitle=\Large\bfseries,
	before=\par\medskip\noindent,parbox=false,boxsep=0pt,left=0pt,right=3mm,top=4pt,
	breakable,pad at break*=0mm,vfill before first,
	overlay unbroken={\draw[colexam,line width=1pt]
		([yshift=-1pt]title.north east)--([xshift=-0.5pt,yshift=-1pt]title.north-|frame.east)
		--([xshift=-0.5pt]frame.south east)--(frame.south west); },
	overlay first={\draw[colexam,line width=1pt]
		([yshift=-1pt]title.north east)--([xshift=-0.5pt,yshift=-1pt]title.north-|frame.east)
		--([xshift=-0.5pt]frame.south east); },
	overlay middle={\draw[colexam,line width=1pt] ([xshift=-0.5pt]frame.north east)
		--([xshift=-0.5pt]frame.south east); },
	overlay last={\draw[colexam,line width=1pt] ([xshift=-0.5pt]frame.north east)
		--([xshift=-0.5pt]frame.south east)--(frame.south west);},%
}

\DeclareTotalTCBox{\commandbox}{ s v }
{verbatim,colupper=pink,colback=black!75!white,colframe=black}
{\IfBooleanT{#1}{\textcolor{cyan}{\ttfamily\bfseries SELECT }}%
	\lstinline[language=command.com,keywordstyle=\color{blue!35!white}\bfseries]^#2^}

\DeclareTotalTCBox{\commandboxa}{ s v }
{verbatim,colupper=white,colback=gray!75!white,colframe=gray}
{\IfBooleanT{#1}{\textcolor{cyan}{\ttfamily\bfseries SELECT }}%
	\lstinline[language=command.com,keywordstyle=\color{blue!35!white}\bfseries]^#2^}

\begin{document}
	
	El aplicar el teorema del elemento extra (EET) a circuitos de primer orden implica los siguientes pasos:
	
	\begin{enumerate}
		\item Identifique el elemento extra $Z$. Puede ser un elemento de almacenamiento de energía $L$ o $C$ o inlcuso una resistencia $R$. El teorema del elemento extra (EET) también funciona para fuentes dependientes. El elemento extra usualmente es seleccionado como el elemento que nos "molesta" y con el cual la función de transferencia se vuelve más compleja de resolver.
		
		\item Decida si cortocircuitar o remover el elemento extra. En algunos casos, si se remueve el elemento la función de transferencia puede ir a cero y no se podría aplicar la ecuación 3.52. Este es el caso para circuitos que tienen un cero en el origen por ejemplo. En su lugar, cortocircuite el elemento y use la ecuación 3.55. Una vez que el elemento extra es colocado en su \textit{estado de referencia} (abierto o cortocircuitado) calcula el término principal $A|_{z=0}$ o $A|_{z=\infty}$. Este término es llamado \textit{ganancia de referencia}.
		
		\item Aplique las ténicas vistas en los capítulos anteriores. Establezca la fuente de excitación a 0 y evalue la resistencia vista desde el puerto creado cuando el elemento extra es removido. con esto se obtendrá $Z_d$
		
		\item Use la doble inyección null (Null Double Injection NDI) para determinar la resistencia en el puerto creado cuando el elemento extra es removido y la respuesta es nula. Con esto se obtendrá $Z_n$
		
		\item Si el circuito de referencia es puramente resistivo, $Z_d = R_d$ y $Z_n = R_n$ son resistencias. El factor de correción directamente nos da directamente las frecuencia de esquina del circuito bajo estudio.
		
	\end{enumerate}

	\begingroup
	\Large
	\begin{align}
		A|_Z = A|_{Z=\infty} \frac{1 + \frac{Z_n}{Z} }{1 + \frac{Z_d}{Z} }
	\end{align}
	\endgroup
	\begingroup
	\Large
	\begin{align}
		A|_Z = A|_{Z=0} \frac{1 + \frac{Z}{Z_n} }{1 + \frac{Z}{Z_d} }
	\end{align}
	\endgroup
	
	
\end{document} 

